\section{Introduction}
\label{submission:intro}

The proliferation of task-specific deep neural networks, fine-tuned from large-scale pre-trained models, has created a compelling need for efficient multi-task integration. Traditional multi-task learning requires expensive joint optimization, which is often computationally prohibitive and susceptible to catastrophic forgetting. Consequently, \emph{model merging}~\cite{izmailov2018swa, wortsman2022soups} has emerged as an attractive, gradient-free alternative. By directly interpolating the weights of pre-trained models fine-tuned on separate tasks, model merging enables the creation of a unified multi-task model with zero retraining cost.

Early approaches to model merging, such as Task Arithmetic~\cite{ilharco2022taskvectors}, rely on static, uniform coefficients to blend task-specific weight deviations (task vectors) with a base pre-trained model. While computationally trivial, static merging suffers from task interference and fails to adapt to input-dependent variations. To overcome this limitation, recent research has pivoted toward \emph{dynamic model merging}~\cite{lu2023adamerging}, which adjusts merging coefficients at test-time based on the input representations. These dynamic methods employ router networks to output sample-specific coefficients, allowing the model to adaptively route inputs to the most relevant task experts.

Despite their empirical promise, existing dynamic model merging techniques are characterized by a profound lack of mathematical rigor. Many state-of-the-art approaches rely on heuristic formulations or complex physical metaphors, such as wavefunctions and quantum eigenstates~\cite{qwsmerge}, without establishing a formal learning-theoretic foundation to explain \emph{why} or \emph{when} they generalize. In this work, we argue that this theoretical deficit has blinded the field to two critical, systemic vulnerabilities:

\begin{figure}[t]
\vskip 0.2in
\begin{center}
\centerline{\includegraphics[width=\columnwidth]{results_plot.png}}
\caption{Classification accuracy (\%) of R2D-Merge and baseline methods under homogeneous and batch-averaged heterogeneous (collapsed) streams. While unregularized linear routers and SOTA quantum-inspired models suffer catastrophic performance collapse when coefficients are averaged across the batch dimension, R2D-Merge demonstrates absolute resilience (0.00\% drop), outperforming the next-best unregularized method by 11.50\% and QWS-Merge by 5.50\% in the collapsed state.}
\label{fig:accuracy_collapse}
\end{center}
\vskip -0.2in
\end{figure}

\begin{enumerate}
    \item \textbf{Transductive Overfitting:} Dynamic routers are optimized at test-time or calibrated on extremely sparse datasets (e.g., $N=64$ samples). In this low-data regime, unconstrained high-dimensional routing functions easily overfit to transductive stream noise or local calibration features. Consequently, their performance on out-of-distribution (OOD) tasks collapses, violating the fundamental goal of general multi-task capability.
    \item \textbf{Heterogeneity Collapse:} In realistic edge deployments, deep neural networks process batches containing samples from different tasks (heterogeneous streams). However, executing sample-specific merged models on standard hardware requires independent forward passes for each sample, scaling the computational cost to $O(B)$ where $B$ is the batch size. To preserve the $O(1)$ single-model efficiency, hardware engines average the dynamic coefficients over the batch dimension: $\bar{\alpha} = \frac{1}{B} \sum_{b=1}^B \alpha_b$. This averaging collapses the task-specific coefficients back to a uniform average, destroying the specialized parameter configurations and causing a catastrophic drop in classification accuracy.
\end{enumerate}

To address these vulnerabilities, we approach dynamic model merging from first principles. Rather than inventing another over-parameterized heuristic, we analyze the statistical generalization of dynamic routing. By formally bounding the empirical Rademacher complexity of the dynamic parameter-space blending function class, we prove that the generalization error is bounded by a weighted Frobenius norm of the router weights. 

This theoretical insight leads directly to our proposed framework: \textbf{Rademacher-Regularized Dynamic Model Merging (R2D-Merge)}. R2D-Merge comprises three core components:
\begin{itemize}
    \item \textbf{Low-Dimensional Projection:} We project input representations into a highly compressed, low-dimensional space ($d = 4$) using a frozen, unsupervised Principal Component Analysis (PCA) matrix, followed by unit-sphere normalization. This restricts the representation capacity of the router.
    \item \textbf{Layer-Wise Linear Routing:} We employ parameter-efficient linear routers that map the compressed input state directly to layer-specific merging coefficients, eliminating redundant wave-like parameters.
    \item \textbf{Covariance-Weighted Frobenius Regularization (CFR):} We introduce a novel quadratic regularizer derived from our Rademacher complexity bound. CFR weights the penalty on the router weights by the task-specific feature covariance matrices scaled by the localized energy of the task vectors. Because these covariance matrices are computed offline during calibration, CFR introduces \emph{zero online computational overhead}.
\end{itemize}

Empirically, we evaluate R2D-Merge on a Vision Transformer (ViT-Tiny) backbone across four diverse image classification datasets: MNIST, FashionMNIST, CIFAR-10, and SVHN. As shown in Figure~\ref{fig:accuracy_collapse}, R2D-Merge achieves comparable accuracy on homogeneous streams while exhibiting absolute resilience to heterogeneity collapse (0.00\% accuracy drop), outperforming unregularized global routers by 11.50\% and quantum-inspired methods by 5.50\% in the collapsed state.

\textbf{Contributions.} Our main contributions are summarized as follows:
\begin{enumerate}
    \item We provide the first formal learning-theoretic generalization bound for dynamic model merging, using empirical Rademacher complexity to analyze parameter-space blending.
    \item We derive Covariance-weighted Frobenius Regularization (CFR), a task-adaptive regularizer that directly minimizes the generalization error bound and is pre-computed offline with zero inference overhead.
    \item We empirically demonstrate that R2D-Merge completely eliminates heterogeneity collapse in multi-task edge streams, providing a mathematically sound and practically robust solution for dynamic model merging.
\end{enumerate}
